% !TEX root = main.tex
% ==================================================================
%  POSTER A0 -- Camara Termica Compacta com Controle PWM
%  Seminario Integrador / Infor -- Engenharia de Computacao -- UNIP 2026
%
%  COMO EDITAR (grupo):
%   - Este arquivo so tem o layout e o cabecalho (titulo, autores, logos).
%   - O texto de cada secao fica em secoes/*.tex  -> cada um edita a sua.
%   - Imagens: coloque em figuras/ (prefira .jpg; .png so para graficos/prints).
%     Enquanto a imagem nao existir, aparece um quadro cinza "Inserir figura".
%   - Tudo que esta em VERMELHO (\pendente{...}) ainda falta preencher.
%   - Plano gratuito do Overleaf: compilacao maxima de 10 s. Imagens grandes
%     deixam lento -> reduza as fotos para ~2000 px de largura antes de subir.
% ==================================================================
\RequirePackage{fix-cm}               % fontes da matematica escalaveis (tamanhos grandes)
\documentclass[a0,portrait]{a0poster}

\usepackage[utf8]{inputenc}
\usepackage[T1]{fontenc}
\usepackage[brazil]{babel}
\usepackage[scaled]{helvet}            % mesma fonte do poster de referencia
\renewcommand{\familydefault}{\sfdefault}
\usepackage{courier}                   % e-mail em fonte monoespacada
\usepackage[margin=2.5cm, bottom=3.5cm]{geometry}
\usepackage{xcolor}
\usepackage{graphicx}
\usepackage{multicol}
\usepackage{booktabs}
\usepackage{amsmath}
\DeclareFontShape{OMX}{cmex}{m}{n}{<->cmex10}{}   % evita aviso de tamanho da fonte matematica
\usepackage{enumitem}
\usepackage{tikz}
\usetikzlibrary{arrows.meta, positioning, calc}
\usepackage{eso-pic}                   % rodape fixo no pe da pagina
\usepackage[hidelinks]{hyperref}

% ---------------- Cores e estilos ----------------
\definecolor{azultitulo}{RGB}{20,40,190}     % azul dos titulos (igual a referencia)
\definecolor{cinzabox}{RGB}{238,238,238}
\setlength{\columnsep}{3.5cm}
\setlength{\parindent}{1.5em}
\setlength{\parskip}{0.3ex}
\setlist{leftmargin=1.2em, itemsep=0.2ex, topsep=0.4ex}

% Titulo de secao azul centralizado:  \secao{Motivacao}
\newcommand{\secao}[1]{\par\vspace{1.4cm}%
  {\centering\color{azultitulo}\LARGE\bfseries #1\par}\vspace{0.6cm}}

% Titulo de figura/tabela (negrito, maiusculo, centralizado, ACIMA da imagem)
\newcommand{\titulofig}[1]{\par\vspace{0.6cm}{\centering\small\bfseries\MakeUppercase{#1}\par}\vspace{0.4cm}}

% Texto pendente (em vermelho) -- some quando voce substitui pelo texto final
\newcommand{\pendente}[1]{{\color{red}[#1]}}

% \imagem[largura][altura]{figuras/arquivo.jpg}
%   cabe a imagem dentro de largura x altura sem distorcer.
%   Se o arquivo nao existir, desenha um quadro cinza com o nome.
\NewDocumentCommand{\imagem}{O{\linewidth} O{12cm} m}{%
  \begin{center}
  \IfFileExists{#3}%
    {\includegraphics[width=#1,height=#2,keepaspectratio]{#3}}%
    {\fcolorbox{gray}{cinzabox}{\parbox[c][\dimexpr#2-2\fboxsep-2\fboxrule][c]{\dimexpr#1-2\fboxsep-2\fboxrule}%
      {\centering\color{gray}\itshape Inserir figura:\\ \normalfont\ttfamily\small\detokenize{#3}}}}%
  \end{center}}

% Logos do cabecalho (quadro cinza se nao existirem)
\newcommand{\logo}[2]{\IfFileExists{#2}%
  {\includegraphics[width=#1,height=#1,keepaspectratio]{#2}}%
  {\fcolorbox{gray}{cinzabox}{\parbox[c][#1][c]{#1}{\centering\color{gray}\small logo\\\ttfamily\detokenize{#2}}}}}

\pagestyle{empty}

% ---------------- Rodape fixo (canto inferior esquerdo) ----------------
\AddToShipoutPictureFG*{%
  \AtPageLowerLeft{\put(\LenToUnit{2.5cm},\LenToUnit{1.6cm}){%
    \small\textbf{Seminário Integrador / Infor 2026}, Engenharia de Computação -- UNIP,
    Novembro de 2026, São José dos Campos, Brasil}}}

\begin{document}

% ================================================================
%  CABECALHO
% ================================================================
\noindent
\begin{minipage}[c]{0.16\textwidth}\centering
  \logo{11cm}{logos/unip.png}
\end{minipage}%
\begin{minipage}[c]{0.68\textwidth}\centering
  {\color{azultitulo}\veryHuge\bfseries
   CÂMARA TÉRMICA COMPACTA COM\\ CONTROLE PWM E CONTROLADOR PI\par}
  \vspace{1.2cm}
  {\Large\bfseries Gabriel Marques\par}
  \vspace{0.7cm}
  {\large Engenharia de Computação -- 8\textsuperscript{o} semestre\\
   Universidade Paulista -- UNIP\\
   São José dos Campos - SP, Brasil\par}
  \vspace{0.5cm}
  {\large\ttfamily gabrielmarques2005@hotmail.com\par}
\end{minipage}%
\begin{minipage}[c]{0.16\textwidth}\centering
  \logo{11cm}{logos/curso.png}
\end{minipage}

\vspace{1cm}

% ================================================================
%  CORPO (2 colunas)
% ================================================================
\begin{multicols}{2}
  % !TEX root = ../main.tex
% ================== MOTIVACAO + INTRODUCAO ==================
\secao{Motivação}

\begin{itemize}
  \item Sistemas de controle de temperatura em malha fechada são a base de fornos,
        incubadoras, estufas e diversos processos industriais, mas plataformas didáticas
        para estudá-los costumam ser caras, volumosas ou de difícil instrumentação.
  \item Propor uma bancada compacta, de baixo custo e totalmente instrumentada para observar
        e validar, de forma visual e mensurável, um controlador PI atuando sobre um sistema
        térmico real.
\end{itemize}

\secao{Introdução}

O controle de temperatura em malha fechada combina instrumentação, eletrônica de potência e
teoria de controle. Este trabalho apresenta uma câmara térmica compacta
(\(\approx\)100\,×\,100\,×\,80\,mm) aquecida por três lâmpadas automotivas W5W, com sensor
LM35 e um controlador PI embarcado em um ESP32, que ajusta o PWM aplicado às lâmpadas por meio
de um MOSFET IRLZ44N para manter a temperatura interna em um setpoint definido [1, 5].

\titulofig{Objetivos específicos}
\begin{itemize}
  \item Estudar controle PI, PWM e instrumentação com sensores analógicos;
  \item Projetar e montar a câmara (acrílico, lâmpadas W5W, chapa de alumínio, LM35 e cadeia
        de potência protegida);
  \item Implementar o firmware do ESP32 (ADC, PWM via LEDC, PI com anti-windup e compensação
        de bateria);
  \item Realizar o ensaio de degrau, sintonizar o PI e validar a malha fechada.
\end{itemize}

  % !TEX root = ../main.tex
% ================== SOLUCAO PROPOSTA ==================
\secao{Solução proposta}

A câmara é uma caixa de acrílico com três lâmpadas W5W (12\,V/5\,W) no teto, uma chapa de
alumínio como massa térmica e o sensor LM35 suspenso no centro. O ESP32 lê a temperatura,
calcula a ação de controle PI e gera o PWM que chaveia as lâmpadas pelo MOSFET.

% ---------- Diagrama da malha de controle (desenhado em TikZ, nao precisa de imagem) ----------
\titulofig{Malha de controle da câmara térmica}
\begin{center}
\begin{tikzpicture}[
    font=\normalsize,
    bloco/.style={draw, thick, rounded corners=6pt, fill=azultitulo!8, align=center,
                  minimum height=3.8cm, minimum width=8.2cm},
    seta/.style={-{Latex[length=7mm]}, line width=2.5pt},
    node distance=2.2cm]
  \node (sp)   {Setpoint};
  \node[draw, thick, circle, minimum size=1.6cm, right=1.6cm of sp] (soma) {};
  \node at (soma) {\(\Sigma\)};
  \node[bloco, right=of soma] (pi)  {Controlador PI\\(anti-windup)};
  \node[bloco, right=of pi]   (pwm) {PWM LEDC\\500\,Hz / 10 bits};
  \node[bloco, below=2.6cm of pwm] (mos) {MOSFET IRLZ44N\\+ 3 lâmpadas W5W};
  \node[bloco, left=of mos]   (cam) {Câmara\\(acrílico + alumínio)};
  \node[bloco, left=of cam]   (sen) {LM35 + ADC1\\(média de 16)};
  \draw[seta] (sp) -- (soma);
  \draw[seta] (soma) -- node[above]{\(e\)} (pi);
  \draw[seta] (pi) -- node[above]{\(u\)} (pwm);
  \draw[seta] (pwm) -- (mos);
  \draw[seta] (mos) -- (cam);
  \draw[seta] (cam) -- node[above]{\(T\)} (sen);
  \draw[seta] (sen.west) -| node[pos=0.8, left]{\(-\)} (soma.south);
\end{tikzpicture}
\end{center}

\vspace{0.5cm}
O controlador PI calcula o duty cycle a partir do erro \(e = T_{sp} - T\):
\begin{equation*}
  u(t) = K_p\, e(t) + K_i \int e(t)\, dt
\end{equation*}
Como a potência nas lâmpadas cai com o quadrado da tensão, o duty é corrigido pela tensão da
bateria:
\begin{equation*}
  D_{corr} = D_{PI} \cdot \left(\frac{12}{V_{bat}}\right)^{2}
\end{equation*}

% ---------- Cadeia de potencia (TikZ) ----------
\titulofig{Cadeia de potência e proteções}
\begin{center}
\begin{tikzpicture}[
    font=\normalsize,
    b/.style={draw, thick, rounded corners=6pt, fill=cinzabox, align=center,
              minimum height=3.4cm, minimum width=6.6cm},
    seta/.style={-{Latex[length=6mm]}, line width=2pt},
    node distance=1.3cm]
  \node[b] (bat) {Pack 3S 18650\\+ BMS};
  \node[b, right=of bat] (d)   {Schottky\\SS34};
  \node[b, right=of d]   (f)   {Fusível\\3\,A + TVS};
  \node[b, right=of f]   (k)   {Termostato\\KSD9700};
  \node[b, below=1.6cm of k] (l) {Lâmpadas\\W5W};
  \node[b, below=1.6cm of d] (buck) {Buck\\12\,\(\to\)\,5\,V};
  \node[b, left=of buck] (esp) {ESP32\\DevKit v1};
  \draw[seta] (bat) -- (d);
  \draw[seta] (d) -- (f);
  \draw[seta] (f) -- (k);
  \draw[seta] (k) -- (l);
  \draw[seta] (d) -- (buck);
  \draw[seta] (buck) -- (esp);
\end{tikzpicture}
\end{center}

\titulofig{Protótipo da câmara térmica}
\imagem[0.85\linewidth][15cm]{figuras/prototipo.jpg}


  \columnbreak   % <- tudo abaixo vai para a coluna da direita

  % !TEX root = ../main.tex
% ================== METODOLOGIA ==================
\secao{Metodologia}

O desenvolvimento seguiu o fluxo \textbf{pesquisa bibliográfica \(\to\) concepção \(\to\)
implementação \(\to\) testes \(\to\) avaliação dos resultados}.

\begin{itemize}
  \item \textbf{Leitura:} LM35 (10\,mV/°C) no ADC1 do ESP32 (12 bits), média de 16 amostras
        a cada 60\,s;
  \item \textbf{Atuação:} PWM via LEDC (500\,Hz, 10 bits), mantido ativo em \emph{light sleep}
        para economizar bateria;
  \item \textbf{Identificação:} ensaio de degrau em malha aberta para extrair a constante de
        tempo \(\tau\) da planta;
  \item \textbf{Sintonia e validação:} ajuste de \(K_p\) e \(K_i\) e análise da resposta ao
        setpoint (overshoot, tempo de acomodação e erro em regime);
  \item \textbf{Segurança:} proteções térmica (KSD9700), elétrica (fusível, TVS, Schottky) e de
        firmware (limite de temperatura).
\end{itemize}

Ferramentas: Arduino IDE / PlatformIO (C++), multímetro, fonte de bancada e osciloscópio.

  % !TEX root = ../main.tex
% ================== RESULTADOS ==================
% Troque os quadros cinza pelos graficos reais (figuras/*.png) e preencha a tabela.
\secao{Resultados e Discussão}

\titulofig{Ensaio de degrau em malha aberta}
\imagem[0.9\linewidth][12.5cm]{figuras/degrau-malha-aberta.png}

\pendente{Descrever o ensaio: duty aplicado, temperatura inicial e final, e o valor de \(\tau\)
obtido (tempo para atingir 63,2\,\% da variação).}

\titulofig{Resposta em malha fechada ao setpoint}
\imagem[0.9\linewidth][12.5cm]{figuras/malha-fechada.png}

\titulofig{Desempenho do controlador PI}
\begin{center}
\begin{tabular}{lc}
  \toprule
  \textbf{Parâmetro} & \textbf{Valor}\\
  \midrule
  Constante de tempo \(\tau\)        & \pendente{-- s}\\
  Ganhos \(K_p\) / \(K_i\)           & \pendente{-- / --}\\
  Setpoint                           & \pendente{-- °C}\\
  Overshoot                          & \pendente{-- \%}\\
  Tempo de acomodação                & \pendente{-- s}\\
  Erro em regime                     & \pendente{\(\pm\) -- °C}\\
  \bottomrule
\end{tabular}
\end{center}

\pendente{Discutir os resultados: o PI manteve a temperatura no setpoint? Houve overshoot? Como
a compensação de bateria afetou a resposta?}

  % !TEX root = ../main.tex
% ================== CONSIDERACOES FINAIS + AGRADECIMENTOS ==================
\secao{Considerações Finais}

O projeto integra sistemas embarcados, eletrônica de potência, instrumentação e teoria de
controle em uma bancada compacta e replicável. As técnicas empregadas -- PWM de baixa
frequência, leitura analógica filtrada, compensação por tensão de bateria e múltiplas camadas
de segurança -- são diretamente transferíveis a produtos industriais de controle térmico.
\pendente{Completar com a conclusão a partir dos resultados obtidos.}

\secao{Agradecimentos}

\pendente{Os autores agradecem ao(à) professor(a) ... pela orientação e à Universidade Paulista
(UNIP) pela infraestrutura disponibilizada.}

  % !TEX root = ../main.tex
% ================== REFERENCIAS ==================
\secao{Referências}

{\small
\begin{enumerate}[label={[\arabic*]}, leftmargin=2.2em, itemsep=0.4ex]
  \item OGATA, K. \emph{Engenharia de Controle Moderno}. 5.\ ed. São Paulo: Pearson, 2010.
  \item ESPRESSIF SYSTEMS. \emph{ESP32 Series Datasheet}. Disponível em:
        https://www.espressif.com, 2024.
  \item TEXAS INSTRUMENTS. \emph{LM35 Precision Centigrade Temperature Sensors}. Datasheet.
        Disponível em: https://www.ti.com/product/LM35, 2017.
  \item VISHAY SILICONIX. \emph{IRLZ44N Power MOSFET}. Datasheet. Disponível em:
        https://www.vishay.com, 2020.
  \item ÅSTRÖM, K. J.; HÄGGLUND, T. \emph{PID Controllers: Theory, Design, and Tuning}.
        2.\ ed. ISA, 1995.
\end{enumerate}
}

\end{multicols}

\end{document}
